\documentclass[12pt, a4paper]{article}

% --- Pacotes Fundamentais ---
\usepackage[utf8]{inputenc}
\usepackage[T1]{fontenc}
\usepackage[brazil]{babel}
\usepackage{geometry}
\usepackage{titlesec}
\usepackage{fancyhdr}
\usepackage{graphicx}
\usepackage{xcolor}
\usepackage{booktabs}
\usepackage{tabularx}
\usepackage{colortbl}
\usepackage{float}
\usepackage{parskip}
\usepackage{lmodern}
\usepackage{lastpage}
\usepackage{amsmath}
\usepackage{hyperref}

% Configuração de hyperlinks
\hypersetup{
    colorlinks=true,
    linkcolor=primaryColor,
    urlcolor=accentColor,
    pdftitle={Memorial de Cálculo - Custo Operacional de Colaborador},
    pdfauthor={Departamento Administrativo}
}

% --- Configurações de Layout ---
\geometry{left=2.5cm, right=2.5cm, top=4.5cm, bottom=2cm, headsep=1.5cm, headheight=60pt}

% --- Definição de Cores ---
\definecolor{primaryColor}{RGB}{0, 51, 102}     % Azul Marinho
\definecolor{secondaryText}{RGB}{80, 80, 80}    % Cinza Escuro
\definecolor{accentColor}{RGB}{178, 34, 34}     % Vermelho Tijolo
\definecolor{lightGray}{RGB}{245, 245, 245}
\definecolor{tableHeader}{RGB}{230, 230, 250}
\definecolor{successGreen}{RGB}{34, 139, 34}

% --- Configuração de Cabeçalho e Rodapé ---
\pagestyle{fancy}
\fancyhf{}

% Cabeçalho
\fancyhead[C]{%
    \begin{tabularx}{\textwidth}{@{} X r @{}}
        \textcolor{primaryColor}{\textbf{\footnotesize MEMORIAL DE CÁLCULO: CUSTO OPERACIONAL DE COLABORADOR}} & \scriptsize \textit{Dept. Administrativo} \\
        \textcolor{secondaryText}{\scriptsize Studio Dental Comércio de Higiene Oral Ltda.} & {\scriptsize \textbf{Ref.: Janeiro/2026}}
    \end{tabularx}%
}

% Rodapé
\fancyfoot[L]{\scriptsize \textcolor{secondaryText}{Documento de uso interno | \textbf{Studio Dental}}}
\fancyfoot[R]{\normalsize Página \textbf{\thepage} de \textbf{\pageref{LastPage}}}

% --- Linha separadora do CABEÇALHO ---
\renewcommand{\headrulewidth}{1.5pt}
\renewcommand{\headrule}{%
    \vspace{0.1cm}
    \hbox to\headwidth{\color{primaryColor}\leaders\hrule height \headrulewidth\hfill}%
}

% --- Linha separadora do RODAPÉ ---
\renewcommand{\footrulewidth}{1.5pt}
\renewcommand{\footrule}{%
    \color{primaryColor}%
    \hbox to\headwidth{\leaders\hrule height \footrulewidth\hfill}%
    \vspace{0.2cm}
}

% --- Formatação de Títulos ---
\titleformat{\section}
{\color{primaryColor}\normalfont\Large\bfseries}
{\thesection}{1em}{}[{\titlerule[0.8pt]}]

\titleformat{\subsection}
{\color{primaryColor}\normalfont\large\bfseries}
{\thesubsection}{1em}{}

% --- Macro para Moeda ---
\newcommand{\reais}[1]{R\$ #1}

% --- INÍCIO DO DOCUMENTO ---
\begin{document}

% --- CAPA ---
\begin{titlepage}
    \centering
    \vspace*{3cm}
    {\Huge \textbf{\textcolor{primaryColor}{MEMORIAL DE CÁLCULO}}} \\
    \vspace{0.5cm}
    {\Large \textbf{CUSTO OPERACIONAL DE COLABORADOR}} \\
    \vspace{0.3cm}
    {\large Análise Detalhada de Custos para Cessão de Mão de Obra} \\
    \vspace{2cm}
    \rule{10cm}{1.5pt} \\
    \vspace{0.5cm}
    \colorbox{lightGray}{%
        \parbox{0.9\textwidth}{%
            \centering
            \vspace{0.3cm}
            {\Large \textbf{GABRYELLA DOS SANTOS SILVA}} \\
            \vspace{0.2cm}
            {\large Auxiliar em Saúde Bucal} \\
            \vspace{0.3cm}
        }%
    }
    \vspace{2cm}

    \begin{flushleft}
    \textbf{COMPOSIÇÃO DO CUSTO:} \\
    \vspace{0.2cm}
    \textcolor{primaryColor}{\textbf{1}} - Remuneração Mensal Direta \\
    \textcolor{primaryColor}{\textbf{2}} - Encargos Sociais e Provisões \\
    \textcolor{primaryColor}{\textbf{3}} - Benefícios Mensais \\
    \textcolor{primaryColor}{\textbf{4}} - Custo por Período e Valor de Repasse
    \end{flushleft}

    \vfill
    {\small Documento elaborado pelo Departamento Administrativo para fins de apuração do custo real da colaboradora, visando a definição de valor justo para cessão/empréstimo de mão de obra entre profissionais ou clínicas parceiras.} \\
    \vspace{1.5cm}

    \colorbox{accentColor}{%
        \parbox{0.9\textwidth}{%
            \centering
            \vspace{0.6cm}
            \textcolor{white}{\Large \textbf{STUDIO DENTAL}} \\
            \vspace{0.2cm}
            \textcolor{white}{\normalsize Comércio de Higiene Oral Ltda.} \\
            \vspace{0.6cm}
        }%
    }

    \vspace{0.2cm}
    \textbf{Data de Referência:} Janeiro de 2026
\end{titlepage}

% --- CONTEÚDO ---

\section{Identificação}

\begin{table}[H]
    \centering
    \renewcommand{\arraystretch}{1.5}
    \begin{tabularx}{\textwidth}{X X}
        \toprule
        \rowcolor{tableHeader} \textbf{Campo} & \textbf{Informação} \\
        \midrule
        Empresa & Studio Dental Comércio de Higiene Oral Ltda. \\
        \rowcolor{lightGray}
        Colaboradora & Gabryella dos Santos Silva \\
        Cargo & Auxiliar em Saúde Bucal \\
        \rowcolor{lightGray}
        Data de Referência & Janeiro/2026 \\
        \bottomrule
    \end{tabularx}
\end{table}

\section{Remuneração Mensal Direta}

Valores extraídos do demonstrativo de pagamento, considerando o saldo base integral para fins de cálculo de custo.

\begin{table}[H]
    \centering
    \renewcommand{\arraystretch}{1.5}
    \begin{tabularx}{\textwidth}{X r}
        \toprule
        \rowcolor{tableHeader} \textbf{Descrição} & \textbf{Valor (R\$)} \\
        \midrule
        Salário Base (Saldo Base) & 1.591,84 \\
        \rowcolor{lightGray}
        Insalubridade & 270,90 \\
        Assiduidade (5\%) & 68,98 \\
        \rowcolor{lightGray}
        Ajuda de Custo & 454,00 \\
        \midrule
        \rowcolor{accentColor!10} \textbf{Total Remuneração Mensal} & \textbf{2.385,72} \\
        \bottomrule
    \end{tabularx}
\end{table}

\section{Encargos Sociais e Provisões}

Cálculo dos custos ``invisíveis'' --- provisões legais e encargos patronais estimados para garantir a cobertura de férias, 13º salário e demais obrigações trabalhistas.

\begin{table}[H]
    \centering
    \renewcommand{\arraystretch}{1.5}
    \begin{tabularx}{\textwidth}{X c r}
        \toprule
        \rowcolor{tableHeader} \textbf{Descrição} & \textbf{Base/Alíquota} & \textbf{Valor (R\$)} \\
        \midrule
        FGTS & 8\% & 154,53 \\
        \rowcolor{lightGray}
        Provisão de 13º Salário & 1/12 & 160,91 \\
        Provisão de Férias + 1/3 & 11,11\% & 214,60 \\
        \rowcolor{lightGray}
        Encargos Patronais (INSS/RAT)* & 28\% & 540,88 \\
        \midrule
        \rowcolor{accentColor!10} \textbf{Total de Encargos/Provisões} & & \textbf{1.070,92} \\
        \bottomrule
    \end{tabularx}
\end{table}

\textcolor{secondaryText}{\footnotesize{*Nota: O valor de INSS Patronal pode variar conforme o enquadramento tributário (Simples Nacional). Alíquota estimada para fins de planejamento.}}

\section{Benefícios Mensais}

Considerando a premissa de \textbf{22 dias úteis} por mês para cálculo de vale transporte.

\begin{table}[H]
    \centering
    \renewcommand{\arraystretch}{1.5}
    \begin{tabularx}{\textwidth}{X c r}
        \toprule
        \rowcolor{tableHeader} \textbf{Descrição} & \textbf{Cálculo} & \textbf{Valor (R\$)} \\
        \midrule
        Vale Alimentação (VA) & Média Mensal & 300,00 \\
        \rowcolor{lightGray}
        Vale Transporte (VT) & R\$ 8,60/dia $\times$ 22 dias & 189,20 \\
        \midrule
        \rowcolor{accentColor!10} \textbf{Total de Benefícios} & & \textbf{489,20} \\
        \bottomrule
    \end{tabularx}
\end{table}

\section{Resumo do Custo Total}

Consolidação de todos os componentes do custo operacional mensal da colaboradora.

\begin{table}[H]
    \centering
    \renewcommand{\arraystretch}{1.6}
    \begin{tabularx}{\textwidth}{X r}
        \toprule
        \rowcolor{tableHeader} \textbf{Componente} & \textbf{Valor Mensal (R\$)} \\
        \midrule
        Remuneração Direta & 2.385,72 \\
        \rowcolor{lightGray}
        Encargos Sociais e Provisões & 1.070,92 \\
        Benefícios Mensais & 489,20 \\
        \midrule
        \rowcolor{primaryColor!15} \textbf{CUSTO TOTAL MENSAL} & \textbf{3.945,84} \\
        \bottomrule
    \end{tabularx}
\end{table}

\section{Cálculo do Custo por Período}

Para definir o valor de cessão/empréstimo da colaboradora, dividimos o custo total pelo número de dias úteis e, posteriormente, pelo número de períodos disponíveis por dia.

\begin{table}[H]
    \centering
    \renewcommand{\arraystretch}{1.5}
    \begin{tabularx}{\textwidth}{X r}
        \toprule
        \rowcolor{tableHeader} \textbf{Parâmetro} & \textbf{Valor} \\
        \midrule
        Custo Total Mensal & \reais{3.945,84} \\
        \rowcolor{lightGray}
        Dias Úteis Considerados & 22 dias (seg. a sex.) \\
        Períodos por Dia & 2 (manhã e tarde) \\
        \bottomrule
    \end{tabularx}
\end{table}

\vspace{0.5cm}

\colorbox{lightGray}{%
    \parbox{\textwidth}{%
        \vspace{0.5cm}
        \centering
        {\large \textbf{Fórmulas de Rateio}} \\
        \vspace{0.5cm}
        \[ \text{Custo por Dia} = \frac{\reais{3.945,84}}{22} = \reais{179,36} \]
        \vspace{0.3cm}
        \[ \boxed{\textbf{Custo por Período} = \frac{\reais{179,36}}{2} = \mathbf{\reais{89,68}}} \]
        \vspace{0.3cm}
    }%
}

\section{Considerações para Cobrança}

O valor de \textbf{\reais{89,68}} representa o \textbf{custo real} da colaboradora para a empresa por período trabalhado. Para fins de cessão ou empréstimo de mão de obra entre profissionais ou clínicas, recomenda-se:

\begin{table}[H]
    \centering
    \renewcommand{\arraystretch}{1.5}
    \begin{tabularx}{\textwidth}{X c r}
        \toprule
        \rowcolor{tableHeader} \textbf{Cenário} & \textbf{Margem} & \textbf{Valor por Período (R\$)} \\
        \midrule
        Custo Real (base) & 0\% & 89,68 \\
        \rowcolor{lightGray}
        Com Margem de Risco Mínima & 15\% & 103,13 \\
        \textcolor{successGreen}{\textbf{Valor Sugerido para Repasse}} & \textcolor{successGreen}{\textbf{20\%}} & \textcolor{successGreen}{\textbf{107,62}} \\
        \bottomrule
    \end{tabularx}
\end{table}

\newpage
\textcolor{accentColor}{\textbf{Recomendação:}} Adotar margem de \textbf{15\% a 20\%} sobre o custo real, resultando em valor de repasse entre \textbf{\reais{105,00} e \reais{110,00}} por período, cobrindo:

\begin{itemize}
    \item Riscos trabalhistas e eventuais passivos
    \item Custos administrativos de gestão
    \item Desgaste de equipamentos e insumos
    \item Margem de segurança para variações sazonais
\end{itemize}

\vspace{2cm}

\begin{center}
    \rule{8cm}{0.5pt} \\
    \vspace{0.3cm}
    \textbf{Responsável pelo Cálculo} \\
    \textcolor{secondaryText}{\small Departamento Administrativo --- Studio Dental}
\end{center}

\end{document}
